\section{Introduction}\label{sec:introduction}

% An introduction where you explain the aims and rationale for the physics case
% and what you have done. At the end of the introduction you should give a brief
% summary of the structure of the report

% Introduction (~1 page). Motivation (why polynomials, why Runge, the
% high-degree failure, bias–variance), then what you did for linear models,
% resampling and optimizers, stating explicitly what is your own code and what
% is scikit-learn. It ends with a paragraph on the report's structure and the
% repo link. No results, no MSE/

Polynomial regression is among the simplest nontrivial regression problems, yet
it already exhibits the central tension in model selection: increasing the
degree reduces the approximation error, but increases sensitivity to noise in
the data. The Weierstrass approximation theorem~\cite{weierstrass_1885} guarantees that any continuous
function on a closed interval can be approximated uniformly to arbitrary
accuracy by polynomials, but it says nothing about how to find such a polynomial
from a finite set of noisy samples, or which degree to choose. Increasing the
degree reduces the approximation error of the model class, yet it also makes the
fitted coefficients increasingly sensitive to the particular data set. This
tension is the bias--variance trade-off~\cite{hjorth-jensen_2026, hastie_2017}. Its lessons carry over essentially
unchanged to the more flexible models used elsewhere in machine learning, which
makes polynomial regression a transparent laboratory for them.

Runge's function,
\begin{equation}\label{eq:runges_function}
  f(x) = \frac{1}{1+25x^2}, \qquad x\in[-1,1],
\end{equation}
is the classical example of how this can fail. Although $f$ is smooth,
its polynomial interpolants on equispaced nodes diverge near the endpoints
as the degree increases;
this is Runge's phenomenon~\cite{runge_1901,demarchi_2020}. %verify runge1901
Least-squares regression on noisy data is not interpolation,
but it faces the same underlying difficulty in a statistical form.
A high-degree fit on equispaced points is free to oscillate between the data,
it fits the noise as readily as the signal,
and the monomial design matrix becomes severely ill-conditioned.
Runge's function therefore provides a demanding test of how model complexity
should be selected and controlled.

We study three linear models. Ordinary least squares (OLS) provides the
unregularized baseline,
whose test error as a function of degree exposes the trade-off directly.
Ridge regression~\cite{hoerl_1970} and Lasso regression~\cite{tibshirani_1996} %verify
add an $\ell_2$ and an $\ell_1$ penalty, respectively, trading a controlled increase in bias for
a reduction in variance. The Ridge penalty shrinks all coefficients, while the Lasso penalty can
set some of them exactly to zero. To estimate the test error and select the degree and penalty,
we use two resampling methods: the bootstrap, which we also use to decompose the error into bias
and variance, and $k$-fold cross-validation.

OLS and Ridge have closed-form solutions, whereas Lasso does not, since its
penalty is not differentiable at zero. We therefore also solve all three problems
iteratively: OLS and Ridge with plain gradient descent, momentum, AdaGrad,
RMSprop and Adam~\cite{kingma_2017}, full-batch and with
mini-batches (stochastic gradient descent), and Lasso with the same update rules
applied to a subgradient of its cost. The closed-form
OLS and Ridge solutions serve as exact benchmarks for these optimizers. We
compute gradients both analytically and by automatic
differentiation~\cite{baydin_2018}, and verify that the two agree.

All regression models, the gradient-descent solvers, the analytical
gradients and the bootstrap are our own implementations. Cross-validation uses
the \texttt{KFold} and \texttt{cross\_validate} utilities of
\texttt{scikit-learn}~\cite{pedregosa_2011} around our
estimators, automatic differentiation uses
\texttt{JAX}~\cite{bradbury_2018}, and
\texttt{scikit-learn}'s \texttt{Lasso} serves as an independent check of our
Lasso solver. Because our data are synthetic, we can also compute the true test
error of every fitted model from the known $f$, and we use it to judge how well
the resampling estimates select models.

The remainder of this report is structured as follows. \Cref{sec:method}
presents the models, the bias--variance decomposition, the resampling methods
and the optimization methods. \Cref{sec:implementation} describes the data,
the preprocessing, the code and how it was verified.
\Cref{sec:results-discussion} presents and discusses the results in the order
OLS, Ridge, bias--variance, resampling and model selection, gradient descent,
Lasso, and stochastic gradient descent, and \cref{sec:conclusion} summarizes our
findings and outlines future work. Appendix~\ref{app:ai-tools} declares our use of
large language models. All code, and the results behind every figure and table,
are available at \url{https://github.com/SigurdRiseth/FYS-STK4155-P1}.
